% ============================================================
\chapter*{Pr\'eface}
\addcontentsline{toc}{chapter}{Pr\'eface}
\markboth{Pr\'eface}{Pr\'eface}
% ============================================================

% ------------------------------------------------------------
\section*{Objectifs de ce cours}
% ------------------------------------------------------------

Le \emph{Traitement Automatique du Langage} (TAL), ou \emph{Natural Language
Processing} (NLP), constitue l'un des domaines les plus dynamiques de
l'intelligence artificielle contemporaine. Ce cours de niveau Master/Doctorat
vise \`a fournir une formation compl\`ete, allant des fondements math\'ematiques
classiques jusqu'aux mod\`eles de langue de derni\`ere g\'en\'eration.

L'objectif principal est triple :
\begin{enumerate}
    \item \textbf{Comprendre les fondements th\'eoriques} --- mod\'elisation
          probabiliste du langage, th\'eorie de l'information, alg\`ebre
          lin\'eaire pour les repr\'esentations vectorielles.
    \item \textbf{Ma\^itriser les architectures modernes} --- r\'eseaux
          r\'ecurrents, m\'ecanismes d'attention, Transformers, mod\`eles
          pr\'e-entra\^in\'es (BERT, GPT, T5).
    \item \textbf{D\'evelopper un esprit critique} --- \'evaluation rigoureuse,
          consid\'erations \'ethiques, biais algorithmiques, s\'ecurit\'e des
          syst\`emes g\'en\'eratifs.
\end{enumerate}

% ------------------------------------------------------------
\section*{Public vis\'e}
% ------------------------------------------------------------

Ce cours s'adresse aux \'etudiants de Master~2 et de Doctorat en informatique,
en science des donn\'ees ou en linguistique computationnelle. Les pr\'erequis
suivants sont suppos\'es :

\begin{itemize}
    \item \textbf{Math\'ematiques} : alg\`ebre lin\'eaire (espaces vectoriels,
          d\'ecompositions matricielles), calcul diff\'erentiel multivari\'e,
          probabilit\'es et statistiques.
    \item \textbf{Informatique} : programmation en Python, notions
          d'algorithmique, familiarit\'e avec les biblioth\`eques scientifiques
          (NumPy, PyTorch ou TensorFlow).
    \item \textbf{Apprentissage automatique} : r\'egression, classification,
          descente de gradient, r\'eseaux de neurones de base.
\end{itemize}

% ------------------------------------------------------------
\section*{Organisation du cours}
% ------------------------------------------------------------

Le cours est organis\'e en onze chapitres, suivant une progression logique
depuis les repr\'esentations classiques du texte jusqu'aux syst\`emes
g\'en\'eratifs les plus r\'ecents :

\begin{description}
    \item[Chapitre~1 --- Introduction au TAL.]
        Panorama historique, t\^aches fondamentales, d\'efis sp\'ecifiques au
        langage naturel.
    \item[Chapitre~2 --- Repr\'esentations classiques du texte.]
        Mod\`eles sac-de-mots, TF-IDF, $n$-grammes, repr\'esentations creuses.
    \item[Chapitre~3 --- Plongements de mots.]
        Word2Vec (CBOW, Skip-gram), GloVe, FastText ; propri\'et\'es
        g\'eom\'etriques des espaces de plongements.
    \item[Chapitre~4 --- Mod\`eles de s\'equences.]
        R\'eseaux r\'ecurrents (RNN), LSTM, GRU ; probl\`emes de gradient
        \'evanescent et explosif.
    \item[Chapitre~5 --- Attention et Transformers.]
        M\'ecanisme d'attention, auto-attention, architecture Transformer
        compl\`ete.
    \item[Chapitre~6 --- Mod\`eles pr\'e-entra\^in\'es.]
        BERT, GPT, T5 ; paradigme pr\'e-entra\^inement/adaptation.
    \item[Chapitre~7 --- Ajustement fin et instruction tuning.]
        Strat\'egies d'adaptation, LoRA, RLHF, alignement.
    \item[Chapitre~8 --- G\'en\'eration de texte et d\'ecodage.]
        D\'ecodage glouton, recherche en faisceau, \'echantillonnage, top-$k$,
        top-$p$.
    \item[Chapitre~9 --- G\'en\'eration augment\'ee par r\'ecup\'eration (RAG).]
        Recherche de passages, int\'egration de connaissances externes, pipelines
        RAG.
    \item[Chapitre~10 --- \'Evaluation des grands mod\`eles de langue.]
        M\'etriques automatiques, \'evaluation humaine, benchmarks, robustesse.
    \item[Chapitre~11 --- \'Ethique, biais et s\'ecurit\'e.]
        Biais dans les donn\'ees et les mod\`eles, toxicit\'e, vie priv\'ee,
        r\'egulation.
\end{description}

% ------------------------------------------------------------
\section*{Conventions typographiques}
% ------------------------------------------------------------

\begin{itemize}
    \item Les \textbf{d\'efinitions} formelles sont encadr\'ees dans des
          environnements \texttt{definition}.
    \item Les \textbf{th\'eor\`emes}, \textbf{propositions} et \textbf{lemmes}
          sont num\'erot\'es et, lorsque pertinent, accompagn\'es d'une preuve.
    \item Les blocs \textbf{intuition} fournissent une compr\'ehension
          qualitative avant la formalisation.
    \item Les blocs \textbf{attention} signalent les erreurs fr\'equentes.
    \item Les blocs \textbf{bonne pratique} r\'esument les recommandations
          concr\`etes.
    \item Le code Python utilise principalement les biblioth\`eques
          \texttt{transformers} (Hugging Face), \texttt{torch} et
          \texttt{scikit-learn}.
\end{itemize}

% ------------------------------------------------------------
\section*{Notations math\'ematiques}
% ------------------------------------------------------------

Nous adoptons les notations suivantes tout au long du cours :

\begin{center}
\renewcommand{\arraystretch}{1.3}
\begin{tabular}{cl}
\toprule
\textbf{Notation} & \textbf{Signification} \\
\midrule
$\R^d$ & Espace euclidien de dimension $d$ \\
$\N$ & Ensemble des entiers naturels \\
$\mathbf{x}, \mathbf{W}$ & Vecteurs (minuscule gras), matrices (majuscule gras) \\
$\inner{\mathbf{u}}{\mathbf{v}}$ & Produit scalaire de $\mathbf{u}$ et $\mathbf{v}$ \\
$\norm{\mathbf{x}}$ & Norme euclidienne de $\mathbf{x}$ \\
$\abs{S}$ & Cardinalit\'e d'un ensemble $S$ \\
$\E[X]$ & Esp\'erance de la variable al\'eatoire $X$ \\
$P(A \mid B)$ & Probabilit\'e conditionnelle de $A$ sachant $B$ \\
$\KL(p \| q)$ & Divergence de Kullback-Leibler de $p$ vers $q$ \\
$\sigma(\cdot)$ & Fonction sigmo\"ide logistique \\
$\softmax(\cdot)$ & Fonction softmax \\
$\argmax_x f(x)$ & Argument maximisant $f$ \\
$\argmin_x f(x)$ & Argument minimisant $f$ \\
$\mathcal{V}$ & Vocabulaire (ensemble fini de tokens) \\
$V = \abs{\mathcal{V}}$ & Taille du vocabulaire \\
$\mathbf{E} \in \R^{V \times d}$ & Matrice de plongements \\
$T$ & Longueur d'une s\'equence \\
\bottomrule
\end{tabular}
\end{center}

% ------------------------------------------------------------
\section*{Logiciels et environnement}
% ------------------------------------------------------------

Les travaux pratiques et exemples de code s'appuient sur l'\'ecosyst\`eme
suivant :

\begin{itemize}
    \item \textbf{Python} $\geq 3.10$
    \item \textbf{PyTorch} $\geq 2.0$
    \item \textbf{Hugging Face Transformers} $\geq 4.35$
    \item \textbf{Hugging Face Datasets}
    \item \textbf{scikit-learn} pour les m\'ethodes classiques
    \item \textbf{NLTK} et \textbf{spaCy} pour le pr\'etraitement linguistique
    \item \textbf{Matplotlib} / \textbf{seaborn} pour la visualisation
\end{itemize}

\begin{bonnepratique}[title=\textbf{Environnement reproductible}]
Nous recommandons d'utiliser un environnement virtuel (\texttt{venv} ou
\texttt{conda}) et de fixer les graines al\'eatoires (\texttt{torch.manual\_seed},
\texttt{random.seed}, \texttt{numpy.random.seed}) pour garantir la
reproductibilit\'e des exp\'eriences.
\end{bonnepratique}

% ------------------------------------------------------------
\section*{Perspective historique}
% ------------------------------------------------------------

Le TAL a connu plusieurs r\'evolutions paradigmatiques. Dans les ann\'ees 1950,
les premi\`eres approches \'etaient fond\'ees sur des r\`egles linguistiques
cod\'ees manuellement. Les ann\'ees 1990 ont vu l'essor des m\'ethodes
statistiques (mod\`eles de Markov cach\'es, mod\`eles $n$-grammes). Les
ann\'ees 2010 ont \'et\'e domin\'ees par l'apprentissage profond (plongements
de mots, r\'eseaux r\'ecurrents). Depuis 2017, l'architecture Transformer
a inaugur\'e l'\`ere des grands mod\`eles de langue (LLM), transformant
radicalement le paysage du domaine.

\begin{intuition}[title=\textbf{Pourquoi le langage est-il difficile ?}]
Le langage naturel est intrins\`equement ambigu (polysmie, anaphores,
ellipses), contextuel (le sens d\'epend du contexte pragmatique) et
compositional (le sens d'une phrase d\'epend de la structure syntaxique et
du sens de ses constituants). Ces propri\'et\'es rendent le TAL
fondamentalement diff\'erent du traitement de donn\'ees structur\'ees.
\end{intuition}

% ------------------------------------------------------------
\section*{Comment utiliser ce cours}
% ------------------------------------------------------------

Chaque chapitre est con\c{c}u pour \^etre auto-suffisant, tout en s'appuyant
sur les concepts des chapitres pr\'ec\'edents. Nous recommandons une lecture
s\'equentielle pour une premi\`ere approche, puis un usage cibl\'e comme
r\'ef\'erence.

Les exercices propos\'es \`a la fin de chaque chapitre sont class\'es en trois
niveaux de difficult\'e :
\begin{itemize}
    \item[$\star$] Exercices de compr\'ehension et de v\'erification.
    \item[$\star\star$] Exercices d'approfondissement n\'ecessitant une
          r\'eflexion math\'ematique ou une impl\'ementation.
    \item[$\star\star\star$] Exercices de recherche ouverts, pouvant servir de
          point de d\'epart pour un projet ou un m\'emoire.
\end{itemize}

% ------------------------------------------------------------
\section*{Remerciements}
% ------------------------------------------------------------

Ce cours a b\'en\'efici\'e des contributions de nombreux coll\`egues,
chercheurs et \'etudiants. Nous remercions en particulier les communaut\'es
open source de Hugging Face, PyTorch et spaCy, dont les outils rendent
l'enseignement du TAL \`a la fois rigoureux et accessible.

Nous remercions \'egalement les \'etudiants des promotions pr\'ec\'edentes
dont les questions et retours ont grandement am\'elior\'e la qualit\'e de ce
mat\'eriel p\'edagogique.

\vfill
\begin{flushright}
\textit{[Author Name], mars 2026}
\end{flushright}
